\section{Adjacent Pell descent and the eleven-positive-variable equation}
\label{sec:positive11}

The second contribution uses a different coherent tuple, the compact
source-44 numerals with $\beta=2^{590069}-1$. The same letters
$K_0,M_C,M_F$ now refer to that tuple. They are bound separately in
\texttt{coefficients/ACTIVE\_BINDING.json}. The four actual integer
payloads and their complete regeneration recipe are included.

\subsection{A rational coordinate is already integral}

The following descent is source 69's key lemma. Its application here
removes a factor previously inserted solely to prove integrality.

\begin{lemma}[Adjacent-coefficient integrality]
\label{lem:descent}
Let $Q>1$ be an integer. If $V\in\Z$, $y\in\Q$, and
\[
 QV^2-(Q-1)y^2=1,
\]
then $y\in\Z$.
\end{lemma}
\begin{proof}
Changing signs permits $V,y>0$; neither is zero. Write $y=b/m$ in
lowest terms with $m>0$. The equation implies $m^2\mid Q-1$, hence
$m\mid Q-1$. If $V=1$ then $y=1$ and $m=1$.

For $V>1$ define
\[
 V'=(2Q-1)V-2(Q-1)y,\qquad
 y'=(2Q-1)y-2QV.
\]
The first is an integer. The identity
$(2Q-1)^2-4Q(Q-1)=1$ preserves the equation. Its original form gives
$y>V$ and $QV>(Q-1)y$; the latter also follows by squaring positive
quantities, since the difference of their squares is $QV^2+Q-1>0$.
These inequalities give $-V<V'<V$. The preserved norm excludes $V'=0$.
Moreover $y'$ has the same reduced denominator $m$ as $y$ because
$2Q-1\equiv1\pmod m$. Replace both new coordinates by their absolute
values and use induction on the positive integer $V$. The descent reaches
$V=1$ without changing $m$, proving $m=1$.
\end{proof}

The adjacent coefficients and right side $1$ matter. For example,
$2\cdot1^2-4(1/2)^2=1$ has a nonintegral rational coordinate, so no
general statement for arbitrary positive coefficients is being used.

\subsection{The exact eleven-variable polynomial}

Supply precisely
\[
 j,F,\alpha,i,T,X,Y,n,\eta,\zeta,\gamma\in\Z_{>0}.
\]
All symbols in the following displays are polynomial abbreviations:
\begin{align*}
 J&=3j+1,&q&=\beta J+1,&C&=q-F-\alpha,\\
 E&=XY,&k&=\eta+\zeta,&c&=kY+\eta,\\
 a&=E+Y,&H&=4a+3,&\Delta&=a^2+H,\\
 \ell&=EkY,&I&=ic^2,&D_\tau&=\ell(\ell+k)+1.
\end{align*}
The removal of $q(q-1)$ from $I$ is the structural change from source 70.
Define
\begin{align*}
 R&=\bigl((q-1)(q-F)+\alpha\bigr)(q^2-1)+(M_C+qM_F)J,\\
 \widehat R&=q^3(q+F+\alpha),& M&=q^3(q-1),\\
 B_t&=(qK_0+X)C-qF,&D_f&=1+\Delta I^2,\\
 Q&=\Delta(D_f-1)=\Delta^2I^2,&\delta&=Q-1,\\
 V&=D_f(cT-R)-c,&\overline V&=D_f(cT+\widehat R),\\
 L&=3\overline V,&W&=1+(q-1)L,\\
 \overline Z&=M(\ell+k+n+L)+(q-1)Y-q^3WB_t.
\end{align*}
For the retained main carrier put
\begin{align*}
 v_0&=k-R-1,& S&=4qY,&U&=qS,\\
 D_*&=U(X+ac)+H(U\gamma+SX-v_0),\\
 A_1&=D_*^2-U^2(1+\Delta c^2).
\end{align*}
For the three-root norm set
\begin{align*}
 A&=M^2D_\tau,&B_f&=M^2D_f,&C_n&=M^2(QV^2-1),\\
 T_n&=\delta^2AB_f,&R_n&=\delta C_n,\\
 J_n&=\delta(\overline Z^2-A-B_f)-C_n,\\
 N_c&=J_n^2-4\bigl(T_n+R_n(A+B_f)\bigr),\\
 N_8&=N_c^2-64T_nR_n\overline Z^2.
\end{align*}
The polynomial is
\begin{equation}
 \boxed{P_{11}=N_8+A_1+N_8A_1=0.}
 \label{eq:P11}
\end{equation}
This literal division-free circuit has 111 binary operations.

\subsection{Radical recovery without a prime-support multiplier}

The universal identity
\begin{align*}
 F_8(z;a,b,c)
 &=\bigl((z^2-a-b-c)^2-4(ab+ac+bc)\bigr)^2-64abc z^2\\
 &=\prod_{\varepsilon\in\{\pm1\}^3}
   (z-\varepsilon_1\sqrt a-\varepsilon_2\sqrt b-\varepsilon_3\sqrt c)
\end{align*}
is a polynomial identity. For positive rational radicands all eight
displayed roots are real; it does not require independent square classes.

We recall source 70's elementary radical lemma. If three positive rational
radicands are pairwise distinct and no signed sum of their square roots
is zero, any rational signed sum of them makes each root rational.
Indeed let signed roots $x_1,x_2,x_3$ have nonzero rational sum $z$,
and write $a_i=x_i^2$. The quantities
\[
 u=\frac{z^2-a_1-a_2-a_3}{2},\qquad
 v=\frac{u^2-a_1a_2-a_1a_3-a_2a_3}{2z}
\]
are rational and equal to $\sum_{i<j}x_ix_j$ and $x_1x_2x_3$.
Then
\[
 x_1=\frac{a_1(a_1+u)}{a_1z+v},\qquad
 a_1z+v=x_1(x_1+x_2)(x_1+x_3)\ne0.
\]
The other two formulas are cyclic. Distinct radicands exclude the two
possible cancellations in each denominator.

Here the three radicands before multiplication by $M^2$ are
\[
 D_\tau,\qquad D_f,\qquad D_y=\frac{QV^2-1}{Q-1}.
\]
All of the following bounds hold for arbitrary positive supplied inputs,
before any equation or the inequality $C>0$ has been assumed. In
particular $c\ge3$, $a\ge2$, $\Delta\ge15$, $I\ge c^2$, and
\[
 1<D_\tau<D_f<Q.
\]
Put $b=cT-R$. If $b=0$, $V=-c$; otherwise
$|V|\ge D_f-c>1$. Thus $|V|>1$ in all cases and
$D_y=V^2+(V^2-1)/(Q-1)>V^2$.
If $D_y$ equalled either integer $D\in\{D_\tau,D_f\}$, then
$Q(D-V^2)=D-1$, impossible because $0<D-1<Q$.

Writing $t_\tau=\sqrt{D_\tau}$ and $f=\sqrt{D_f}$ gives
\[
 t_\tau<(a+1)c,\qquad f>(a+1)c^2\ge3c^2.
\]
When $b=0$, $\sqrt{D_y}<c+1<f-t_\tau$. When $b\ne0$,
$\sqrt{D_y}>D_f-c>2f>f+t_\tau$. One root therefore dominates
the sum of the other two, excluding a zero signed sum.

The authentic mask bounds give
$0<M_C<\beta$ and $0<M_F<2\beta$. They imply
$|R|<\widehat R$ for every positive input. For example
\[
 |R|<q^4+Fq^3+(\alpha+2)q^2<q^3(q+F+\alpha),
\]
using $q\ge16$, $\alpha\ge1$. The positive integer gap, together with
$D_f>c$, gives $|V|<\overline V$. Hence
\[
 f<\overline V,\qquad \sqrt{D_y}<2|V|<2\overline V,
 \qquad L>f+\sqrt{D_y},\quad L>k.
\]
These are exactly the estimates needed for source 70's root transport;
none uses the removed factor $q(q-1)$ in $I$.

At a zero of $N_8$, the identity
\[
 N_8=\delta^4 F_8(\overline Z;M^2D_\tau,M^2D_f,M^2D_y)
\]
supplies signed roots with
$\overline Z/M=\tau_*+f_*+y_*$. The radical lemma makes them rational.
The first two are integral because their squares are integers. The third
satisfies
\[
 QV^2-(Q-1)y_*^2=1.
\]
Lemma~\ref{lem:descent} therefore gives $y_*\in\Z$ directly.
This replaces source 70's prime-support argument and is why the stronger
divisibility of $I$ is unnecessary.

\subsection{Soundness and constructive completeness}

For soundness recover
\[
 f_{\rm old}=|f_*|,\quad u_{\rm old}=|y_*|-V,
 \quad r_{\rm old}=n+L-f_*-y_*,\quad i_{\rm old}=i.
\]
The inequalities above make $u_{\rm old},r_{\rm old}$ strictly positive.
The signed roots must be used in $r_{\rm old}$; replacing them there by
absolute values would change the transport identity. With
\[
 Z=M(\ell+k+r_{\rm old})+(q-1)Y-q^3WB_t
\]
one obtains $Z=M\tau_*$. The four source-70 weighted-carrier equations
are now
\begin{align*}
 Z^2-M^2D_\tau&=0,&A_1&=0,\\
 f_{\rm old}^2-\Delta(i_{\rm old}c^2)^2-1&=0,\\
 (V+u_{\rm old})^2-\Delta(f_{\rm old}^2-1)
            u_{\rm old}(2V+u_{\rm old})-1&=0.
\end{align*}
Its theorem applies to every positive integer transport weight, here
$L$. It restores the normalized positive source and hence a finite
refutation.

For clarity, this inherited weighted-carrier theorem has a specific
order of recovery. First $M\mid Z$ yields $q^3\mid Y$ and
$(q-1)\mid B_t$, and the positive transport gap gives $B_t>0$, thus
$C>0$. Only then are the native bounds $3q<R<q^4-2$ used. The adaptive
main carrier gives a positive main Pell root, the strong/auxiliary
equations recover its index, and its large modulus recovers $q\mid X$
and the deleted index quotient. This order avoids using native bounds
before their hypotheses have been established.

Conversely choose canonical normalized source witnesses for a finite
refutation. A positive Pell pair with ordinate divisible by $c^2$ is
available; it follows either from the canonical source or from a positive
power of the Pell matrix that is the identity modulo $c^2$. Thus
$I=ic^2$ with $i>0$. The actual auxiliary reconstruction lemma of source
61 applies to every positive such pair with $c\mid I$, canonical
$R\ge3$, $R\equiv3\pmod4$, and $c=\psi_{a+2}(R)$, supplying
positive $T,u$ and $y=V+u$. Here $\psi_A(R)$ is the coefficient of
$\sqrt{A^2-1}$ in $(A+\sqrt{A^2-1})^R$.
Choose these before computing $L,W$. Canonical transport has
$t_0=B_t/(q-1)\ge q$, $s_0=Y/q^3<k$, and $\ell<\tau<\ell+k$.
Its positive gap can therefore be set to
\[
 r=Wt_0+\tau-\ell-k-s_0>qW-2k>L.
\]
Now $n=r-L+f+y>0$ and $\overline Z=M(\tau+f+y)$, so $N_8=0$.
The adaptive main-carrier witness is supplied by the same canonical
construction. This proves completeness with no extra strong-index
congruence.

Finally $N_8$ is a square modulo $4$, and so is $A_1$ because $4\mid U$.
The equation $P_{11}=0$ is $(1+N_8)(1+A_1)=1$. The negative-unit case
would require both residuals to be $-2$, impossible modulo $4$.
Thus this single equation is equivalent to their simultaneous vanishing.

\subsection{Exact degree and its coefficient certificate}

The exact degrees are
\[
\begin{array}{c|rrrrrrrr}
 \text{register}& I&D_f&Q&V&L&\overline Z&N_8&A_1\\\hline
 \text{degree}&5&14&18&18&18&24&264&13.
\end{array}
\]
The degree-fourteen terms in $A_1$ cancel; the exact audit checks this
identity rather than trusting naive gate degrees. Writing
$q_0=3\beta j$ and $k_0=\eta+\zeta$, the leading forms satisfy
\[
 \lead_{264}(N_8)=Q_0^4\overline Z_0^8,
\]
where
\begin{align*}
 Q_0&=i^2k_0^4X^4Y^8,\\
 B_{t,0}&=(K_0q_0+X)(q_0-F-\alpha)-q_0F,\\
 \overline Z_0&=-3i^2q_0^7k_0^4X^2Y^6(q_0+F+\alpha)B_{t,0},\\
 \lead_{13}(A_1)&=32q_0^2k_0X^2Y^4
                    \bigl(4q_0^2Y\gamma+q_0^3(q_0-F)\bigr).
\end{align*}
In particular the degree-277 monomial has coefficient
\[
 [j^{60}i^{24}\eta^{49}X^{42}Y^{85}\alpha^{16}\gamma]P_{11}
 =2^7\,3^8(3\beta)^{60}>0.
\]
The degree falls by $48$, or about $14.77\%$, from source 70's 325.
The live arithmetic count falls by two multiplications, from 113 to 111.

\section{A twelve-positive-variable companion of degree 78}

This construction composes source 69's norm merger with source 70's
adaptive main carrier. Supply
\[
 j,F,\alpha,t,T,X,Y,r,\eta,\zeta,n,\gamma\in\Z_{>0}.
\]
Use $J=6j+1$, $q=\beta J+1$, and the same compact tuple and common
registers $C,E,k,c,a,H,\Delta,\ell,R,M,B_t$ as above. Set
\begin{align*}
 W&=(q-1)k+1,\\
 Z&=M(\ell+k+r)+(q-1)Y-q^3WB_t,\\
 A_0&=Z^2-M^2\bigl(\ell(\ell+k)+1\bigr).
\end{align*}
Use the adaptive $S=4qY$, $U=qS$, $D_*$ and $A_1$ from the preceding
section, in place of source 69's scales $q^4,q^5$.
The remaining registers are
\begin{align*}
 f&=t+1,&\rho&=f^2-1,&d&=\Delta c^4,&Q&=\Delta\rho,\\
 V&=f^2(cT-R)-c,\\
 h(z)&=(1-Q)z^2+2(1-Q)Vz+V^2-1,\\
 b_N&=2(1-Q)(n+V),&A_N&=d\,h(n)+(1-Q)\rho,\\
 N&=A_N^2-d\,b_N^2\rho.
\end{align*}
The single equation is
\begin{equation}
 \boxed{P_{12}=(1+A_0)(1+A_1)(1+N)-1=0.}
 \label{eq:P12}
\end{equation}

\subsection{The merger and its exact recovery}

Assume $d,\rho>0$, $Q>1$, $V>1$, $n>0$ and $N=0$.
Then $b_N\ne0$. Define the proof-only rational numbers
\[
 i_*=\frac{A_N}{d b_N},\qquad u_*=n-i_*.
\]
One gets $d i_*^2=\rho$ and $h(u_*)=0$. Writing $y_*=V+u_*$,
the latter equation is $QV^2-(Q-1)y_*^2=1$.
Lemma~\ref{lem:descent} gives $y_*\in\Z$, hence $u_*,i_*\in\Z$.
Recover the positive witnesses $i=|i_*|$ and $u=|y_*|-V$;
the second is positive because $|y_*|>|V|$.
Conversely, if $d i^2=\rho$ and $h(u)=0$, set $n=i+u$. The identity
\[
 N\big|_{\rho=di^2}=d^2h(n-i)h(n+i)
\]
shows $N=0$. This is an exact polynomial identity; it is not an
unjustified rational-root resultant inference.

The required hypotheses precede this recovery. The equations $A_0=A_1=0$
alone give $C>0$, $q^3\mid Y$ and positive first and main Pell indices.
The first index is at least two; otherwise $k=2$, $\eta=\zeta=1$
would contradict the smallest possible main ordinate. Therefore
\[
 c>kY>8XY^3\ge8q^9>2R.
\]
Since $f=t+1\ge2$, this gives
$V=f^2(cT-R)-c>c>1$. The adaptive carrier has the same hypotheses and
rank consequences here as in source 70. In particular $q\equiv2\pmod3$
and $q^3\mid Y$ make its relevant gcd at most $3$, while its remaining
modulus exceeds $R+1$. Thus replacing $q^5,q^4$ by $4q^2Y,4qY$
does not weaken the source reconstruction.

The recovered $i,u$ restore the four positive source residuals, with
transport weight $k$. Conversely, canonical witnesses have
$J\equiv1\pmod6$ because the radix exponent and repunit length are odd;
they supply $t=f-1>0$ and $n=i+u>0$. The first transport and adaptive
main-carrier gaps are positive by the canonical construction. This
proves the existence equivalence with the same native source.

The radix $q$ is even, $4\mid U$, and $b_N$ is even. Hence all three
residuals $A_0,A_1,N$ are squares modulo $4$. Their unit-product
finalizer in \eqref{eq:P12} is equivalent to their separate vanishing:
any factor $-1$ would again require a residual $-2$.

\subsection{Degree and arithmetic cost}

The exact residual degrees are $15,13,50$, so the product has degree 78.
An explicit coefficient surviving fixed-port specialization is
\[
 [j^{18}F^2\eta^{11}X^{12}Y^{23}t^8n^2\gamma\alpha]P_{12}
 =1024(6\beta)^{18}>0.
\]
The new scale computation costs one more live gate than the old scale,
giving 103 operations. This is a degree improvement from 82 to 78,
about $4.88\%$, and an explicit tradeoff rather than domination of the
old circuit in every measure.
